% ------------------------------------------------------------------------------
% SLIDES FOR NLP COLLOQUIUM ON 2020-11-27
% ------------------------------------------------------------------------------ 

% SETUP ------------------------------------------------------------------------

\documentclass[t]{beamer}
\usepackage[utf8]{inputenc}

\setbeamercolor{palette primary}{bg = gray!90,fg = white}
\setbeamercolor{structure}{fg = gray!90} % itemize, enumerate, etc
\setbeamercolor{alerted text}{fg = gray!90}
\newcommand{\highlight}[1]{\colorbox{gray!20}{#1}}
\setbeamertemplate{itemize items}[circle]
\newcommand{\arrowitem}{\item[\textcolor{gray!90}{$\rightarrow$}]}

\title{Topic-Specific Sentiment Analysis for Twitter Data of German MPs}
\author{Asmik Nalmpatian, Lisa Wimmer}
\institute{Statistical Consulting}
\date{2020}

\usepackage{Sweave}
\begin{document}
\Sconcordance{concordance:20201127_nlp_colloqium.tex:20201127_nlp_colloqium.Rnw:%
1 21 1 1 0 121 1}


\frame{\titlepage}

\footnotesize

% ------------------------------------------------------------------------------

\begin{frame}
\frametitle{Project Outline}

\textbf{Goal:} \highlight{topic-specific sentiment analysis} on
tweets by German MPs \\
\smallskip
$\rightarrow$ What sentiments are expressed toward particular topics? \\
  
\begin{figure}
  \raggedright
  \includegraphics[width = 0.8\textwidth]{figures/screenshot_tweet.PNG}
\end{figure}

\textbf{Two-pronged approach:}
\begin{enumerate}
  \footnotesize
  \item Comprehensive \highlight{pipeline}: \\
  data extraction > data processing > sentiment analysis
  \begin{itemize}
    \footnotesize
    \item Robust framework with rather basic methodology
    \item To be implemented in \texttt{R}
  \end{itemize}
  \item \highlight{Advanced} sentiment analysis
  \begin{itemize}
    \footnotesize
    \item Exploration of more complex methods
    \item Most probably in \texttt{Python}
  \end{itemize}
\end{enumerate}

\end{frame}

% ------------------------------------------------------------------------------

\begin{frame}
\frametitle{Key Challenges}

\begin{itemize}
  \footnotesize
  \item \textbf{Large data base of unstructured text}
  \begin{itemize}
    \footnotesize
    \arrowitem High dimensionality
    \arrowitem Run time, memory
  \end{itemize}
  \item \textbf{German language}
  \begin{itemize}
    \footnotesize
    \arrowitem Syntactic complexities
    \arrowitem Less existing research than for English
  \end{itemize}
  \item \textbf{Twitter idiosyncrasies}
  \begin{itemize}
    \footnotesize
    \arrowitem Short document length (140 characters)
    \arrowitem Informal language (plus spelling mistakes)
    \arrowitem Special features (hashtags, emojis, ...)
  \end{itemize}
  \item \textbf{Topic extraction}
  \begin{itemize}
    \footnotesize
    \arrowitem Upstream task where same challenges are present
  \end{itemize}
  \item \textbf{No labels}
  \begin{itemize}
    \footnotesize  
    \arrowitem No means of evaluation with data as-is
  \end{itemize}
\end{itemize}

\end{frame}

% ------------------------------------------------------------------------------

\begin{frame}
\frametitle{Ideas}

\begin{itemize}
  \item \textbf{First \& foremost:} \highlight{labels}
  \item \textbf{Then:} \highlight{classification} in three levels of complexity
  \begin{enumerate}
    \footnotesize
    \item \textbf{Dictionary approach}
    \begin{itemize}
      \item Baseline model
      \item n-grams, bag-of-words assumption
      \item Probably low accuracy
    \end{itemize}
    \item \textbf{Classic ML models} 
    \begin{itemize}
      \item Focus on feature extraction
      \item Tried-and-tested classifiers (RF, SVM, ...)
    \end{itemize}
    \item \textbf{BERT and friends} 
    \begin{itemize}
      \item Black-box, high-complexity approaches
      \item Hope: data in, magic out
    \end{itemize}
  \end{enumerate}
  \item \textbf{Eventually:} we know more about...
  \begin{itemize}
    \footnotesize
    \item ... how far we can get with basic to medium approaches
    \item ... by how much we can boost accuracy with adding complexity
  \end{itemize}
\end{itemize}

\end{frame}

% ------------------------------------------------------------------------------

\end{document}
